\documentclass[11pt,a4paper]{article}
\usepackage[spanish]{babel}
\usepackage[utf8]{inputenc}
\usepackage[T1]{fontenc}
\usepackage{geometry}
\usepackage{hyperref}
\usepackage{listings}
\usepackage{xcolor}

\geometry{margin=2.5cm}

\definecolor{codebg}{HTML}{F5F5F5}
\lstset{
  basicstyle=\ttfamily\small,
  backgroundcolor=\color{codebg},
  frame=single,
  breaklines=true,
  showstringspaces=false
}

\title{Gateway Java e Integracion con Microservicio de Usuarios}
\author{ShopSmart}
\date{\today}

\begin{document}
\maketitle

\section{Resumen}
Este documento describe el microservicio \textbf{ShopSmart Gateway}, su configuracion, rutas, despliegue en Docker y su integracion con el microservicio de usuarios. El gateway centraliza el enrutamiento, expone endpoints comunes y aplica resiliencia con circuit breaker.

\section{Arquitectura}
El gateway se construye con \textbf{Spring Cloud Gateway} y actua como punto de entrada unico para la API. Enruta solicitudes hacia el microservicio de usuarios y expone endpoints de salud mediante Actuator.

\subsection*{Componentes}
\begin{itemize}
  \item \textbf{Gateway}: Spring Cloud Gateway con rutas definidas en \texttt{application.yml}.
  \item \textbf{Usuarios}: microservicio independiente en el puerto 8081.
  \item \textbf{Resilience4j}: circuit breaker con fallback local.
  \item \textbf{Actuator}: endpoints de salud y rutas del gateway.
\end{itemize}

\section{Dependencias principales}
\begin{itemize}
  \item Spring Cloud Gateway
  \item Spring Boot Actuator
  \item Resilience4j (circuit breaker)
\end{itemize}

\section{Configuracion}
La configuracion principal esta en \texttt{src/main/resources/application.yml}. Valores clave:
\begin{itemize}
  \item Puerto: 8080
  \item Nombre de app: \texttt{gateway-service}
  \item CORS global habilitado
  \item Timeouts HTTP: 5s connect, 10s response
  \item Actuator: \texttt{/actuator/health}, \texttt{/actuator/info}, \texttt{/actuator/gateway}
\end{itemize}

\section{Rutas del Gateway}
\begin{itemize}
  \item \texttt{/api/v1/auth/**} \textrightarrow{} \texttt{http://usuarios-service:8081}
  \item \texttt{/api/v1/usuarios/**} \textrightarrow{} \texttt{http://usuarios-service:8081}
  \item \texttt{/api/v1/admin/usuarios/**} \textrightarrow{} \texttt{http://usuarios-service:8081}
  \item \texttt{/swagger/usuarios/**} \textrightarrow{} \texttt{http://usuarios-service:8081}
\end{itemize}

\section{Resiliencia y fallback}
Se configura un circuit breaker llamado \texttt{usuariosCircuitBreaker}. En caso de fallo, el gateway responde con el endpoint local:
\begin{lstlisting}
GET /fallback/usuarios
\end{lstlisting}

\section{Ejecucion local}
Desde la carpeta del gateway:
\begin{lstlisting}
mvn spring-boot:run
\end{lstlisting}

\section{Docker}
\subsection*{Dockerfile}
El Dockerfile usa multi-stage build para generar el jar y ejecutar el gateway en JRE 17.

\subsection*{docker-compose}
El gateway se conecta a la red \texttt{shopsmart-net} y monta el \texttt{application.yml} como configuracion externa.

Ejecucion:
\begin{lstlisting}
docker compose up -d --build
\end{lstlisting}

\section{Pruebas rapidas}
\subsection*{Login via Gateway}
\begin{lstlisting}
POST http://localhost:8080/api/v1/auth/login
Content-Type: application/json

{"email":"maria.gonzalez@email.com","password":"Segura123"}
\end{lstlisting}

\subsection*{Rutas cargadas}
\begin{lstlisting}
GET http://localhost:8080/actuator/gateway/routes
\end{lstlisting}

\section{Troubleshooting}
\begin{itemize}
  \item \textbf{404 en rutas}: verificar que el gateway haya cargado las rutas en \texttt{/actuator/gateway/routes}.
  \item \textbf{Fallback activo}: revisar conectividad con \texttt{usuarios-service} y red Docker.
  \item \textbf{CORS}: ajustar en \texttt{application.yml} si hay bloqueos en frontend.
  \item \textbf{Config externa}: confirmar variable \texttt{SPRING\_CONFIG\_ADDITIONAL\_LOCATION=/app/config/}.
\end{itemize}

\end{document}
